\section{引言}

空中-水面-水下任务将无人机（UAV）、无人水面艇（USV）和无人水下航行器（UUV）结合起来，用于观测、检查、搜索和环境监测。它们在机动性与感知能力上的互补，使任务在无线电与声学链路、网关、处理服务和管理主体之间形成依赖关系\citep{li2023survey,wei20233u,pang2026cooperative}。无人机及海洋系统综述将机动性、有限能源、间歇通信、定位不确定性和恢复困难列为运行约束\citep{hayat2016survey,shakhatreh2019unmanned,bae2023survey,gonzlezgarca2020autonomous}。任务决策必须综合在这些不同条件下生成的证据。

这形成了\emph{任务信任缺口}。认证确立主体的身份；内容哈希绑定字节；复制共识则在某种故障模型下建立一致性。依赖证据作出的决策还需要证明：观测在物理上可信、时效性足够，并且适合其预期后果。带签名的声呐报告可能不准确，真实图像可能已经过时，排序正确的记录也可能来自被俘获的设备。无人机和水下安全研究描述了相应的欺骗、干扰、俘获、路由攻击和数据操纵威胁\citep{mekdad2023survey,wang2025survey,hadi2023comprehensive,jiang2019securing,yisa2021security}。将这些控制措施与数据产品的最终用途联系起来，需要明确的主张语义。

区块链方案涉及集群认证、数据采集、成员资格、溯源和审计\citep{han2022identity,tan2022blockchain,karmakar2024blockchain,islam2019bus}。无人机与海事应用综述描述了协调与共享的机会\citep{hafeez2023blockchain,alladi2020applications,mehta2020blockchain,farah2024survey,li2021survey}。其价值取决于哪些参与方必须共同治理状态、必须容忍哪些故障，以及经过认证的数据库或独立见证日志是否已经能够满足这些要求。区块链预言机文献直接指出了外部输入问题：对已记录断言达成一致，并不能证明这些断言与物理事件相符\citep{caldarelli2020oracle,albreiki2020oracles}。

相邻领域的机制提供了互补证据。零信任提供面向特定资源的授权\citep{rose2020zero,syed2022zero}；远程证明区分设备证据、评估和依赖决策\citep{birkholz2023remote}；溯源记录生产者和转换过程\citep{w3c2013provdm,w3c2013provo}；信任推断则对声誉、不确定性和行为变化建模\citep{ahmed2019trust,aaqib2023iot,jsang2016subjective}。这些机制使用不同的词汇，评估不同的主张。它们的组合是本综述的核心议题。

所提出的\OAL{}框架，将实体、平台/设备、链路/路径和数据产品，与四个保障层及六个生命周期阶段联系起来。带类型的依赖图与情境-时间-图（CTG）循环，在将该结构应用于决策时保留共同原因、事件时间和不确定性。OAL 提供主张的组织方式；已有的推断、溯源、策略和共识机制提供实现手段。

本综述作出三项相互关联的贡献。第一，将 OAL 应用于 39 项经过全文检查的机制研究，比较互补机制所提供的证据，并检验覆盖情况对纳入选择和编码选择的敏感性。第二，通过分区-撤销-协调案例，将框架与明确的许可账本模式、可执行的策略回放和通信预算相连接。具有策略验证能力的见证日志基线揭示了两种架构产生等价决策的情形。第三，基于需求的共识选择和可证伪的研究命题，将综合结果转化为评估议程。

民用任务和公开记录的联盟场景构成本文的应用背景。
